%=========================================================================
%  Annals of Nuclear Energy (Elsevier) manuscript
%  Class: elsarticle  -- available by default on Overleaf.
%  Recommended submission build: \documentclass[preprint,review,12pt]{elsarticle}
%  (double spaced, line numbers, figures/tables at end) for first submission.
%  Use [preprint,12pt] for a readable single-spaced PDF while drafting.
%
%  Bibliography: refs.bib, BibTeX, elsarticle-harv (author-year) style.
%  \todocite{...} marks a claim that still needs a source. It renders in bold
%  in the draft PDF -- grep for "CITE:" and clear them all before submission.
%=========================================================================
%=========================================================================
%  HIGHLIGHTS -- TODO after results. Elsevier requires these as a SEPARATE
%  file at submission: 3-5 bullets, max 85 characters each including spaces.
%
%  * TODO
%  * TODO
%  * TODO
%=========================================================================
\documentclass[preprint,12pt]{elsarticle}

\usepackage{amsmath,amssymb}
\usepackage{graphicx}
\usepackage{booktabs}
\usepackage{multirow}
\usepackage{array}
\usepackage[section]{placeins}
\usepackage{lineno}
\linenumbers

% --- convenience macros for the recurring neutronics symbols -------------
\newcommand{\keff}{k_{\mathrm{eff}}}
\newcommand{\Sig}[1]{\Sigma_{\mathrm{#1}}}
\newcommand{\nuSigf}{\nu\Sigma_{\mathrm{f}}}
\newcommand{\Bz}{B_z^{2}}
\newcommand{\MWdkg}{\mathrm{MWd/kgHM}}

% visible marker for a claim that still needs a source
\newcommand{\todocite}[1]{\textbf{[CITE: #1]}}

\journal{Annals of Nuclear Energy}

\begin{document}

\begin{frontmatter}

\title{Physics-Informed Graph Neural Operators for Heterogeneous Two-Dimensional
Reactor Neutronics}

% TODO: author list. One \author per person in final order, each keyed to an
% \affiliation; \corref{cor1} on the corresponding author only.
%
% \author[aff1]{Given Name Surname}
% \ead{email@institution.edu}
% \cortext[cor1]{Corresponding author.}
%
% \affiliation[aff1]{organization={Department},
%                    city={City}, state={ST}, postcode={00000}, country={USA}}

%-------------------------------------------------------------------------
\begin{abstract}
% TODO: write after results are in. Max ~250 words, no references,
% no undefined abbreviations. Structure: problem, approach, benchmarks,
% key quantitative findings, significance.
\end{abstract}

\begin{keyword}
Neutron diffusion \sep Graph neural operator \sep Physics-informed machine
learning \sep Operator learning \sep Digital twin \sep Sodium fast reactor \sep
Fluoride-salt-cooled high-temperature reactor
\end{keyword}

\end{frontmatter}

%-------------------------------------------------------------------------
\section*{Nomenclature}
\addcontentsline{toc}{section}{Nomenclature}
\begin{tabbing}
\hspace{2.6cm}\= \kill
$A$ \> neutron loss operator (leakage, removal, in-scatter), cm$^{-1}$ \\
$F$ \> neutron production (fission) operator, cm$^{-1}$ \\
$\phi_g$ \> scalar neutron flux in group $g$, cm$^{-2}$\,s$^{-1}$ \\
$\keff$ \> effective multiplication factor, -- \\
$D_g$ \> diffusion coefficient in group $g$, cm \\
$\Sig{tr}$ \> transport cross section, cm$^{-1}$ \\
$\Sig{a}$ \> absorption cross section, cm$^{-1}$ \\
$\Sig{r}$ \> removal cross section, cm$^{-1}$ \\
$\Sigma_{\mathrm{s},g'\rightarrow g}$ \> group transfer cross section, cm$^{-1}$ \\
$\nuSigf$ \> fission neutron production cross section, cm$^{-1}$ \\
$\chi_g$ \> fission spectrum fraction in group $g$, -- \\
$\Bz$ \> axial geometric buckling, cm$^{-2}$ \\
$E_{\mathrm{f}}$ \> energy released per fission, J \\
$f$ \> control/shutdown element insertion fraction, -- \\
$\beta$ \> boundary albedo, $J^- = \beta J^+$, -- \\
$\alpha$ \> Marshak partial-current Robin coefficient, -- \\
$s$ \> hexagonal submesh refinement level, -- \\
$G$ \> number of energy groups, -- \\
$N,\;T,\;B,\;E$ \> numbers of mesh nodes, triangles, boundary edges, graph edges, -- \\
$V_i$ \> lumped nodal volume of node $i$, cm$^{3}$ \\
$w_{ig}$ \> volume-integrated group flux weight of constituent $i$, cm\,s$^{-1}$ \\
$\mathcal{R}$ \> discrete PDE residual, -- \\
$\mathcal{L}$ \> training loss, -- \\
\end{tabbing}

\noindent\textit{Abbreviations:} CE, continuous energy; DeepONet, deep operator
network; FEM, finite element method; FHR, fluoride-salt-cooled
high-temperature reactor; FIMA, fissions per initial metal atom; GNN, graph
neural network; GNO, graph neural operator; HALEU, high-assay low-enriched
uranium; IBC, initial and boundary conditions; MLP, multilayer perceptron;
PDE, partial differential equation; PI-GNO, physics-informed graph neural
operator; PINN, physics-informed neural network; SFR, sodium fast reactor;
V\&V, verification and validation.

%=========================================================================
\section{Introduction}
\label{sec:intro}

%--------------------------------------------------------------
\subsection{Motivation}
\label{sec:motivation}

Global electricity demand continues to rise, as nations pursue widespread
electrification while simultaneously expanding computational infrastructure.
The emergence of large-scale artificial intelligence systems,
high-performance computing facilities, and hyperscale data centers has
introduced energy requirements that are expected to grow substantially over
the coming decades. Continued reliance on fossil fuels presents environmental
and economic challenges, including greenhouse gas emissions, air pollution, and
long-term energy security concerns. Consequently, interest in scalable
low-carbon sources capable of providing reliable baseload power has increased.
Nuclear energy remains one of the few mature technologies capable of producing
large amounts of carbon-free electricity in high capacity factors, making
it a central component of future energy systems.

The growing importance of nuclear power has motivated research into advanced
monitoring, control, and safety technologies for both the operating fleet and
next-generation designs. Among these, digital twins have emerged as a
framework for integrating simulation, sensor data, and predictive analytics
into a unified representation of a physical system \citep{kobayashi2024}. By
providing real-time
estimates of reactor state and near-term system behavior, digital twins can
improve situational awareness and support decision-making throughout the
reactor lifecycle. Their effectiveness, however, depends on the availability
of computational models that are both sufficiently accurate and capable of
operating at near real-time speed.

Many of the governing equations encountered in reactor physics and multiphysics
simulation are partial differential equations (PDEs), including neutron
diffusion, neutron transport, heat conduction and thermal-hydraulic balance
equations. Traditional numerical approaches such as the finite difference
method, the finite element method (FEM) and Monte Carlo methods provide highly
accurate solutions, but require substantial computational resources,
particularly when repeated evaluations are needed for uncertainty
quantification, design optimization, or real-time monitoring. Data-driven
surrogates address this cost directly: they still rely on conventional solvers
to generate training data, but evaluate the same problems orders of magnitude
faster once trained, with reported speedups of up to approximately
$4\times10^{3}$ relative to the reference solver
\todocite{source for the $\sim$4000$\times$ speedup figure}.

%--------------------------------------------------------------
\subsection{Related work}
\label{sec:related}

The reactor core is among the most important subsystems to analyze for any
nuclear reactor. As an example of data-driven core analysis, an artificial
neural network has been trained on 10{,}000 randomized core configurations to
learn reactivity evolution and cycle length \citep{palmi2025}. Purely
data-driven methods, however, are susceptible to a loss of physical
consistency, and their black-box architecture is difficult to analyze for
verification and validation (V\&V). Physics-informed modeling offers one route
to reintroducing deterministic structure into such models.

Physics-informed neural networks (PINNs) construct the loss from a standard
mean-squared-error term augmented by the residual of the governing equations,
allowing the PDE and its boundary conditions to be incorporated directly into
the optimization objective \citep{raissi2019}. In neutronics, PINNs have
been applied to the point kinetics equations to predict core power evolution by
embedding the residual of the point kinetics equations together with the
corresponding initial and boundary conditions (IBC) in the loss
\citep{schiassi2022}. Incorporating these residuals encourages the network
to satisfy the governing equations and associated IBC during training, which
improves physical consistency and may support future V\&V workflows.

Whereas PINNs improve physical consistency through PDE-constrained training,
neural operators are designed to learn solution operators for entire families
of PDEs. Unlike conventional networks, which learn mappings between
finite-dimensional vectors, neural operators learn mappings between function
spaces \citep{kovachki2023}. This renders the learned map
discretization-invariant and enables generalization across a family of PDE
parameterizations without training a separate model for each configuration.
Reactor analysis frequently requires repeated evaluation across varying
material distributions, control-rod positions and operating conditions, which
makes reactor neutronics a natural candidate for operator learning. A deep
operator network (DeepONet) has been shown to outperform a fully convolutional
network and a convolutional neural network for two-dimensional neutron
transport in a maze geometry \citep{kobayashi2024}.

While neural operators perform strongly on structured computational domains,
reactor cores are inherently heterogeneous systems composed of fuel
assemblies, control rods, burnable poisons, reflector regions and coolant
channels. The associated neutronics calculations are therefore performed on
irregular meshes, where material interfaces, geometric discontinuities and
localized feedback effects strongly influence the flux distribution. A
representation that naturally encodes geometric connectivity and heterogeneous
material distribution is consequently desirable.

Graph neural operators (GNOs) represent reactor discretizations by treating
mesh cells, material regions or sensor locations as nodes connected through
physically meaningful interactions \citep{li2020}. By encoding spatial
locations, mesh cells or sensor measurements as nodes and their physical
interactions as edges, GNOs learn solution operators through message passing,
and they have demonstrated strong performance on PDEs defined over irregular
geometries and unstructured domains \citep{li2020}. Physics- and
geometry-aware graph neural operators have more recently been trained with
PDE-constrained objectives on both regular and irregular domains
\citep{sarkar2025}. Those developments, however, have been demonstrated on
canonical benchmark PDEs. Application to heterogeneous reactor neutronics --
where the operator must act on a strongly discontinuous material field, respect
material interfaces, and return an eigenvalue in addition to a field -- remains
comparatively unexplored.

%--------------------------------------------------------------
\subsection{Contribution and outline}
\label{sec:contribution}

This work investigates a physics-informed graph neural operator (PI-GNO)
framework for heterogeneous reactor neutronics, combining operator learning,
graph-based geometric representations, and PDE-constrained training to
approximate neutronics solutions on irregular reactor geometries. The specific
contributions are as follows.

\begin{enumerate}
\item A PI-GNO architecture that learns the solution operator of the
      heterogeneous multigroup neutron diffusion eigenvalue problem on
      unstructured reactor meshes, predicting group fluxes and $\keff$ and
      recovering nodal power analytically from the predicted flux rather than
      through an additional learned head.
\item A training objective in which the residual of the assembled discrete
      diffusion eigenvalue problem acts as an explicit physics constraint,
      evaluated on the same finite element operators used to label the data, so
      that the Marshak boundary condition is enforced through the operator
      rather than as a separate penalty.
\item A transport-traceable database for two heterogeneous benchmarks -- a
      sodium-cooled fast reactor and a fluoride-salt-cooled pebble-bed reactor
      -- in which every group constant is collapsed from evaluated nuclear data
      by full-core continuous-energy Monte Carlo with in-situ spectral
      weighting, resolved double heterogeneity, transport-corrected diffusion
      coefficients and burnup, temperature and control-insertion branch cases.
\item A verification harness that solves identical core states with both the
      labeling diffusion model and full-core continuous-energy transport,
      reporting the reactivity bias, control worth and radial power deviation.
\item A controlled benchmark against a Fourier neural operator, a deep operator
      network and a mesh graph network, in which all four models share one
      training loop, one parameter budget, one learning-rate grid and, for the
      accuracy comparison, a purely data-driven objective, so that no result is
      attributable to unequal treatment. The comparison is made along three
      axes -- irregular geometry, discretization transfer and boundary-condition
      transfer -- and each is reported against a measured, model-independent
      error floor: the rasterization loss the grid baseline is forced to incur,
      and the finite element solutions' own mesh dependence.
\end{enumerate}
% NOTE: this list states scope, not findings, so it is safe to keep now. But
% item 4 promises an ablation that has not been run yet -- re-check every item
% against the finished results before submission and cut anything unsupported.

The remainder of this paper is organized as follows. Section~\ref{sec:method}
describes the benchmark problems, the data generation and verification
pipeline, the PI-GNO architecture and the evaluation methodology.
Section~\ref{sec:results} presents the results, and Section~\ref{sec:conclusions}
summarizes the conclusions.

%=========================================================================
\section{Methodology}
\label{sec:method}

%--------------------------------------------------------------
\subsection{General framework}
\label{sec:framework}

This work proposes a physics-informed graph neural operator framework for
approximating the solution operator associated with heterogeneous multigroup
neutron diffusion eigenvalue problems on irregular reactor geometries. The
methodology comprises six stages: (i) definition of benchmark reactor
configurations; (ii) generation of high-fidelity reference solutions with a
finite element diffusion solver; (iii) construction of graph representations
from the discretized reactor meshes; (iv) database generation and
preprocessing; (v) training of the proposed PI-GNO architecture; and
(vi) evaluation against baseline machine learning and neural operator models.
The complete workflow is illustrated in Fig.~\ref{fig:workflow}.

\begin{figure}[htbp]
\centering
% \includegraphics[width=\linewidth]{figures/workflow.pdf}
\framebox[\linewidth][c]{\rule{0pt}{5cm}\textsc{Figure placeholder: end-to-end workflow}}
\caption{End-to-end workflow of the proposed PI-GNO framework, from benchmark
definition and offline nuclear data generation through finite element
reference solutions, graph construction and physics-informed training to
evaluation against baseline architectures.}
\label{fig:workflow}
\end{figure}

Reactor cores are heterogeneous systems composed of fuel assemblies, control
rods, reflector regions, shielding materials, and coolant channels that are not
naturally represented on uniform computational grids. The proposed framework
therefore combines neural operator learning with graph-based representations,
allowing the model to learn the governing diffusion operator while preserving
the geometric connectivity of the heterogeneous core. This combination is
particularly suited to digital twin applications that require both rapid
inference and the ability to operate on complex reactor geometries.

%--------------------------------------------------------------
\subsection{Reactor physics and benchmark problems}
\label{sec:benchmarks}

Two distinct two-dimensional heterogeneous multigroup neutron diffusion
eigenvalue problems are considered, representative of a sodium-cooled fast
reactor and a fluoride-salt-cooled pebble-bed reactor.

The first benchmark is inspired by the Natrium-class sodium fast reactor
(SFR). To represent its neutron spectrum, a configurable number of energy
groups $G$ is employed in which all groups resolve the fast spectrum -- a
high-fast and a slow-fast group at the default $G=2$ -- rather than the
conventional fast--thermal approximation used for light water reactors. The
second benchmark is inspired by the Kairos Power fluoride-salt-cooled
high-temperature reactor (KP-FHR). It also uses a configurable group count,
but resolves a genuinely thermal spectrum, with a fast and a thermal group
separated at a cutoff near 0.625~eV at the default $G=2$. Both reactors pass
through an identical solver, graph-construction and model pipeline; the group
count, material set and feature layout are schema-driven from each sample's
metadata rather than hard coded.

The governing equations are formulated as the coupled generalized eigenvalue
problem
\begin{equation}
A\,\boldsymbol{\phi} \;=\; \frac{1}{\keff}\,F\,\boldsymbol{\phi},
\label{eq:eigen}
\end{equation}
where $A$ is the neutron loss operator incorporating leakage, removal and
in-scattering, $F$ is the fission production operator, $\boldsymbol{\phi}$ is
the group-wise neutron flux and $\keff$ is the effective multiplication factor.
In continuous form, the multigroup diffusion balance for group $g$ reads
\citep{duderstadt1976,stacey2007}
\begin{equation}
-\nabla\!\cdot\!\big(D_g \nabla \phi_g\big)
+ \Sig{r,\mathit{g}}\,\phi_g
= \sum_{g'\neq g} \Sigma_{\mathrm{s},g'\rightarrow g}\,\phi_{g'}
+ \frac{\chi_g}{\keff}\sum_{g'} \nu\Sigma_{\mathrm{f},g'}\,\phi_{g'} .
\label{eq:diffusion}
\end{equation}

The two benchmark configurations comprise heterogeneous material distributions
appropriate to each reactor, summarized in Table~\ref{tab:materials}. The
sodium fast reactor consists of inner and outer fuel, primary and secondary
control assemblies, reflector, shield, homogenized duct (HT9 wall plus sodium
gap) and sodium coolant. The pebble-bed reactor consists of fuel and
graphite-moderator pebbles, control and shutdown elements, graphite reflector,
FLiBe coolant and a steel vessel. Vacuum boundary conditions are imposed on the
outer boundary of both benchmarks. The quantities of interest predicted by the
framework are the group-wise neutron flux distribution, $\keff$, and the nodal
power density computed directly from the predicted flux.

\begin{table}[htbp]
\centering
\caption{Material sets of the two benchmark problems.}
\label{tab:materials}
\begin{tabular}{@{}llc@{}}
\toprule
Benchmark & Material types & Count \\
\midrule
Sodium fast reactor (hex) &
\begin{tabular}[t]{@{}l@{}}
fuel\_inner, fuel\_outer, primary\_control,\\
secondary\_control, reflector, shield, duct, sodium
\end{tabular} & 8 \\[2.2ex]
Pebble-bed reactor (FHR) &
\begin{tabular}[t]{@{}l@{}}
fuel\_pebble, graphite\_pebble, control\_element,\\
shutdown\_element, reflector, FLiBe, vessel
\end{tabular} & 7 \\
\bottomrule
\end{tabular}
\end{table}

To evaluate robustness and generalization, multiple scenarios are considered
for each reactor, including variations in geometry, material layout, reflector
thickness and control- and shutdown-element insertion depth. For the pebble bed
the fuel-to-moderator pebble ratio is additionally varied as a reactivity
lever. These variations enable assessment of the ability of each model to learn
an entire family of neutron diffusion operators rather than individual reactor
solutions.

%--------------------------------------------------------------
\subsection{Data generation}
\label{sec:datagen}

Data generation is split into an offline nuclear data stage and a per-sample
generation stage. The offline stage runs continuous-energy Monte Carlo
transport in OpenMC once per reactor -- unit-cell depletion followed by
full-core branch-case transport -- and caches a self-describing group-constant
table. The per-sample stage builds a meshed core state, looks its cross
sections up in that table, solves the multigroup diffusion eigenvalue problem
with P1 finite elements and emits the graph tensors. The per-sample stage never
calls a transport code, and no hand-tuned cross-section library exists anywhere
in the pipeline: a missing or mis-structured table raises a hard error, so
every constant in every sample is traceable to an evaluated nuclear data
library.

\subsubsection{Geometry generation}
\label{sec:geometry}

Reference solutions are generated for two reactor geometries, both discretized
with P1 finite elements and both passing through the same downstream solver and
graph pipeline. Principal dimensions are summarized in
Table~\ref{tab:geometry}.

\paragraph{Sodium fast reactor (hexagonal lattice)}
The core is constructed on a hierarchical hexagonal-duct lattice at an 18.7~cm
assembly pitch with a 0.37~cm HT9 duct wall and a 0.40~cm inter-assembly sodium
gap. Each assembly is discretized with an identical structured triangular
submesh, establishing spatial invariance across physically identical
assemblies; the outer ring, occupying the outer 18\% of the hexagon radius, is a
homogenized HT9-duct-plus-sodium-gap material, following standard
hexagonal-lattice homogenization practice \citep{scale}. The submesh is
parameterized by a refinement level $s$, the only knob in the pipeline that
changes the discretization without changing the physics and therefore the axis
of the resolution-transfer study of Section~\ref{sec:datafamilies}. Level $s=0$
is a 13-node, 18-triangle fan interior wrapped by an annular ring, and is the
level at which the main datasets are generated; $s\geq1$ builds $R=2s$
concentric rings, giving $1+3R(R+1)$ nodes and $6R^{2}$ triangles per assembly,
that is 19, 61 and 127 nodes per assembly at $s=1,2,3$. Level 0 is a different
mesh topology -- six nodes on its single ring, against $6k$ on ring $k$ of the
general scheme -- so it is held out of the refinement family rather than treated
as its coarsest member, and the level is recorded per sample because the node
count alone does not identify it. Explicit sodium nodes are placed at the
midpoints between neighboring assembly centers, and the inter-assembly gaps
are stitched with a Delaunay triangulation \citep{delaunay1934} over the
duct-ring and gap nodes,
retaining only triangles whose centroids fall outside every hexagon and whose
longest edge is shorter than the lattice pitch, so that the gap is a resolved
physical region rather than a blurred gradient.

Radial roles are assigned by ring index: inner-enrichment fuel,
outer-enrichment fuel, reflector rings, then shield rings, with nine primary
and four secondary control positions -- the docketed Natrium counts
\todocite{Natrium core map: assembly and control-assembly counts} --
deterministically replacing selected fuel positions. Radial enrichment zoning
places the higher enrichment in the outer zone, 19.75~wt\% $^{235}$U against
15.50~wt\% in the inner zone, because the fast-spectrum periphery leaks
strongly and requires additional fissile inventory to flatten the radial power
profile. Reflector- and shield-ring counts and the enrichment-boundary ring
vary per sample. The domain boundary is the set of assembly edges without a
neighboring assembly, which carries the Marshak condition of
Section~\ref{sec:solver}, reducing to vacuum at zero albedo.

\paragraph{Pebble-bed reactor (FHR)}
The core is constructed as a full-diameter cylindrical pebble bed at the
published generic FHR (gFHR) dimensions
\todocite{gFHR dimensional specification / benchmark}, the non-proprietary
KP-FHR surrogate and the only core in that family with an open dimensional
specification. The
core is explicitly not annular: KP-FHR pebbles are buoyant and float upward
through the entire bed cross-section, so no central reflector column is
present. The radial build is a pebble bed to $R=120$~cm, a 60~cm solid graphite
side reflector to $R=180$~cm, then an SS316H barrel (2~cm), FLiBe downcomer
(5~cm) and SS316H reactor vessel (4~cm), carried as a single homogenized vessel
ring to $R=191$~cm. The 309.47~cm active height enters the two-dimensional
model only through the transverse buckling and never as an in-plane dimension.

Pebbles of 2.0~cm radius are placed by random sequential adsorption
\todocite{random sequential adsorption packing} across the
bed and homogenized as one node per pebble; a per-sample fraction of 0.08--0.22
are moderator-only graphite pebbles, which is a physical KP-FHR reactivity
lever because it sets the carbon-to-heavy-metal ratio. Interstitial FLiBe nodes
are placed to reproduce the physical pebble-to-coolant volume split of
approximately 60:40. Ten B$_4$C control elements of 2.6~cm radius, centered
7.9~cm outboard of the bed edge, sit in the side graphite reflector, each a
rigid filled 16-node disk inside a 1.6~cm graphite channel lining, following
the NRC KP-FHR core design and analysis methodology topical report
\citep{kpfhr2024}, in which the reactivity control system
inserts into engineered channels in the side reflector rather than into the
bed. Three X-shaped B$_4$C shutdown elements insert directly into the packed
bed with no graphite thimble, since direct pebble and FLiBe contact is what
gives the reactivity shutdown system its worth, and the pebble packer excludes
their inflated footprints so that pebbles seat against the blades. Element
counts and insertion locations follow the licensed design; element geometry is
taken from the published gFHR rod where the licensed dimensions are
proprietary, and the shutdown-blade size is scaled to the bed and documented as
such. The full node cloud is triangulated with a single Delaunay triangulation
and the boundary is its free-edge set.

Both benchmarks are inspired by publicly available descriptions of the
respective reactors. Dimensions are the published gFHR values for the pebble
bed and representative HALEU metal-fuel SFR values for the fast reactor, whose
vendor core map is not in the public record; where a licensed dimension is
proprietary, a documented non-proprietary surrogate is used and recorded in the
sample metadata. Cross sections, by contrast, are never representative: they
are generated from evaluated nuclear data as described in
Sections~\ref{sec:depletion}--\ref{sec:groupconstants}.

\begin{table}[htbp]
\centering
\caption{Principal geometric and modelling parameters of the two benchmarks.}
\label{tab:geometry}
\begin{tabular}{@{}lll@{}}
\toprule
Parameter & Sodium fast reactor & Pebble-bed reactor \\
\midrule
Lattice type              & Hexagonal duct lattice        & Packed cylindrical bed \\
Assembly pitch            & 18.7 cm                       & -- \\
Duct wall / gap           & 0.37 cm HT9 / 0.40 cm Na      & -- \\
Pebble radius             & --                            & 2.0 cm \\
Bed outer radius          & --                            & 120 cm \\
Side reflector            & Variable ring count           & 60 cm graphite (to $R=180$ cm) \\
Vessel ring               & Shield rings (variable)       & 2 cm barrel + 5 cm FLiBe + 4 cm vessel \\
Outer radius              & Lattice-dependent             & 191 cm \\
Active height             & --                            & 309.47 cm (buckling only) \\
Control elements          & 9 primary, 4 secondary        & 10 B$_4$C in side reflector \\
Shutdown elements         & --                            & 3 X-shaped B$_4$C in bed \\
Enrichment ($^{235}$U)    & 15.50 / 19.75 wt\% (in/out)   & gFHR specification \\
Moderator pebble fraction & --                            & 0.08--0.22 \\
Axial buckling $\Bz$      & $5.4\times10^{-4}$ cm$^{-2}$  & $1.0\times10^{-4}$ cm$^{-2}$ \\
Mesh refinement level $s$ & 0 (main), 1--3 (transfer)     & -- \\
Core size                 & 217--271 assemblies, 3421 nodes & $\approx$1730 pebbles, 7098 nodes \\
Boundary condition        & Marshak, albedo $\beta$       & Marshak vacuum \\
\bottomrule
\end{tabular}
\end{table}

\subsubsection{Unit-cell depletion}
\label{sec:depletion}

Fuel burnup is a dataset axis, so fuel isotopics as a function of burnup are
computed rather than assumed. Following standard lattice-physics practice
\citep{scale}, a representative unit cell is depleted with
\texttt{openmc.deplete} \citep{romano2021} using a coupled operator and a
predictor integrator, and the resulting nuclide
inventories are applied to the full-core spectrum calculation; only the
isotopics originate in the unit cell, while the group-collapse spectrum remains
the full-core one. For the pebble bed the cell is a single fuel pebble with its
explicit TRISO lattice inside a reflective FLiBe cell sized such that the pebble
volume fraction equals the bed packing fraction, so that the pebble sees its
in-bed moderator-to-fuel ratio. For the fast reactor the cell is a single fuel
assembly with an explicit 217-pin lattice (U--10Zr slug, sodium bond, HT9 clad)
inside a reflective HT9 duct.

One cell is depleted per fuel zone. The fast reactor has two enrichment zones
and therefore requires two independent depletion calculations, because an
assembly at 15.50~wt\% reaches a different inventory at a given burnup than one
at 19.75~wt\% -- more plutonium relative to remaining $^{235}$U, and a
different fission-product loading per unit fissile -- and because the depleted
composition overrides the fresh material in the transport model, serving one
zone's isotopics to the other would silently erase the enrichment zoning and
flatten the radial power. The tables are therefore keyed by FEM fuel material,
a missing zone raises an error rather than falling back, and the single-zone
table schema is rejected on load. The pebble bed has a single fuel material and
therefore requires one calculation. Heavy-metal volumes are set analytically
from the resolved geometry, that is from kernel and pin counts, rather than by
stochastic volume calculation.

The production tables cover 0, 2, 20, 40 and 60~$\MWdkg$ for the fast reactor
(ENDF/B-VIII.0 fast chain \citep{brown2018}, 45~W/gHM) and 0, 2, 20, 50, 90,
130, 160 and
190~$\MWdkg$ for the pebble bed (ENDF/B-VIII.0 thermal chain, 30~W/gHM), the
latter reaching the published gFHR peak discharge burnup of approximately 20\%
FIMA. Every grid includes an early point near 2~$\MWdkg$ so that downstream
branch cases interpolate past $^{135}$Xe and $^{149}$Sm equilibrium rather than
across that step change; the sampled burnup ranges carry a floor of
2~$\MWdkg$ for the same reason, since an operating core contains no
xenon-free fuel. Compositions are interpolated linearly in atom density between
depletion points and clamped, never extrapolated, outside the grid.

\subsubsection{Full-core transport group constants}
\label{sec:groupconstants}

Multigroup constants are collapsed from continuous-energy ENDF/B-VIII.0 data
\citep{brown2018} by OpenMC \citep{romano2013} eigenvalue transport on
full-core models that mirror the FEM cores. The
transport models import the role and control assignment of the FEM geometry
modules themselves, so the transport core is the FEM core rather than a
parallel re-derivation of it. Four properties make the resulting constants
defensible.

\begin{description}
\item[In-situ weighting.] Every material is tallied with
\texttt{domain\_type="material"} inside the assembled full core, so that
\begin{equation}
\Sigma_g \;=\; \frac{\displaystyle\int_{E_g} \Sigma(E)\,\phi(E)\,\mathrm{d}E}
                    {\displaystyle\int_{E_g} \phi(E)\,\mathrm{d}E}
\label{eq:collapse}
\end{equation}
uses the spectrum that the material actually experiences. Infinite-medium unit
cells are used nowhere in the collapse: an infinite reflective box of B$_4$C
driven by a Watt source bears no relationship to the spectrum seen by a
reflector-adjacent control element, and its absorption constants are wrong by a
large factor in a thermal system.

\item[Double heterogeneity.] Resonance self-shielding is a geometric effect, so
fuel is resolved explicitly: TRISO particle lattices (UCO kernel, buffer, IPyC,
SiC, OPyC) inside explicit pebbles with a low-density buoyancy core and an
annular fuel shell for the pebble bed, and explicit hexagonal pin lattices
inside each duct for the fast reactor. Constituents share a single universe per
kind, instantiated per pebble or per assembly.

\item[Transport-corrected diffusion.] The diffusion coefficient is taken as
$D = 1/(3\Sig{tr})$ \citep{duderstadt1976} from a tallied transport cross
section, not from a
flux-weighted total cross section.

\item[Branch cases rather than multipliers.] Burnup, temperature and control
insertion are each represented by a genuine transport branch: depleted
isotopics from Section~\ref{sec:depletion}, true Doppler broadening and
$S(\alpha,\beta)$ thermal scattering at the branch temperature, and rodded
versus unrodded cores with the absorber channels filled by either B$_4$C or the
physical follower (sodium for the fast reactor, FLiBe for the pebble bed).
\end{description}

For each branch, the model tallies transport, absorption, fission, fission
neutron production, the $\nu$-scatter matrix -- which carries $(n,2n)$
multiplication, as the diffusion balance requires -- and $\chi$ over the group
structure, together with an explicit per-material group flux used as the
homogenization weight. Group boundaries are defined in exactly one place in the
code and drive the collapse, the fission spectra and the sample metadata alike:
0.1~MeV for the fast reactor, near the flux peak and above the $^{238}$U
inelastic threshold so that the second group captures the slowing-down tail in
which control worth resides, and 0.625~eV for the thermal pebble bed.

The group count $G$ is a configuration parameter rather than a fixed choice, and
a second, finer structure is tabulated for the fast reactor in order to test
whether its two-group collapse is what limits the accuracy of the labels. That
structure is deliberately \emph{nested} in the two-group one: the 0.1~MeV cut is
retained and two further cuts are added at 1.353~MeV and 9.1188~keV, both
lethargy-grid points, so that the four-group indices $\{1,2\}$ sum exactly to the
old fast group and $\{3,4\}$ to the old second group. The added cuts isolate
$^{238}$U threshold fission and high-energy inelastic scattering above
1.353~MeV, where $\nu$ and $\chi$ are hardest, and separate the
unresolved-resonance capture range from the resolved resonances -- sodium at
2.85~keV and the structural iron and chromium -- where the capture-to-fission
ratio climbs most steeply and rod and reflector worth concentrate. Because the
partition is a refinement of the same partition rather than a different one, any
difference between the two is attributable to group resolution alone. The
corresponding four-group fission spectrum is derived from the same Watt spectrum
as the two-group value and collapses consistently with it, and a separate
group-constant table is generated on this structure with the identical full-core
branch procedure. Section~\ref{sec:verification} reports what that comparison
measures. The pebble bed is left at $G=2$ on purpose, since its two-group bias is
already small and needs no refinement.

Resolved constituents are then flux-volume homogenized back onto the single FEM
material carried by each mesh node, following standard assembly-homogenization
practice \citep{scale}. Reaction rates are additive, so the
rate-conserving homogenized constant is
\begin{equation}
\Sigma_{\mathrm{hom},g} \;=\;
\frac{\sum_i w_{ig}\,\Sigma_{ig}}{\sum_i w_{ig}},
\qquad
w_{ig} = \int_{V_i}\!\!\int_{E_g} \phi\,\mathrm{d}E\,\mathrm{d}V ,
\label{eq:homog}
\end{equation}
with $w_{ig}$ the volume-integrated group flux in constituent $i$. The scatter
matrix is weighted by the incoming-group flux, and the diffusion coefficient is
homogenized through its transport cross section,
\begin{equation}
\Sigma_{\mathrm{tr,hom},g} = \frac{\sum_i w_{ig}\,\Sigma_{\mathrm{tr},ig}}{\sum_i w_{ig}},
\qquad
D_{\mathrm{hom},g} = \frac{1}{3\,\Sigma_{\mathrm{tr,hom},g}},
\label{eq:dhomog}
\end{equation}
because averaging $D$ directly is not physical. Because the material-filtered
flux score in OpenMC is already volume-integrated, no further volume factor is
applied; doing so would weight each constituent by $V^2$ and break rate
conservation. Removal is defined as
\begin{equation}
\Sig{r,\mathit{g}} = \Sigma_{\mathrm{a},g} + \sum_{g'\neq g}\Sigma_{\mathrm{s},g\rightarrow g'},
\label{eq:removal}
\end{equation}
that is absorption plus total out-scatter in both directions, and the upper
triangle of the tallied scatter matrix is carried separately as up-scatter data.
A single core-average fission spectrum per branch is formed by weighting the
tallied $\chi$ of each material by its fission neutron production, since $\chi$
is fixed nuclear data belonging in the production operator rather than in the
node features.

\subsubsection{Branch interpolation and cross-section assignment}
\label{sec:branch}

The generator requires cross sections at arbitrary per-node states, whereas
transport can only afford a coarse branch grid, so the table is interpolated in
the standard lattice-physics manner. Burnup and temperature are interpolated
bilinearly between bracketing branch points and clamped, never extrapolated;
macroscopic cross sections are close to linear in both over the tabulated
spacing. A generation-time check warns when a sampled range escapes the
tabulated grid, because outside that grid every state clamps onto the same
point and a slice of the dataset would silently carry identical constants while
still advertising transport provenance.

Control and shutdown insertion is expressed through a depth $f\in[0,1]$, with
both endpoints measured by transport. The cross sections of the element blend
from the tabulated withdrawn follower branch ($f=0$; a sodium follower in the
fast reactor, a FLiBe-filled channel in the pebble bed) toward the tabulated
fully inserted B$_4$C absorber branch ($f=1$). Removal, scattering, and fission
cross sections blend linearly with $f$, corresponding to volume- or
worth-weighted homogenization of a partially present absorber,
\begin{equation}
\Sigma_g(f) = (1-f)\,\Sigma_g^{\mathrm{out}} + f\,\Sigma_g^{\mathrm{in}},
\label{eq:blendsigma}
\end{equation}
while the diffusion coefficient blends linearly in its transport cross section,
\begin{equation}
\frac{1}{3D_g(f)} = (1-f)\,\frac{1}{3D_g^{\mathrm{out}}} + f\,\frac{1}{3D_g^{\mathrm{in}}},
\label{eq:blendd}
\end{equation}
which is equivalent to a worth-weighted harmonic mean of $D$. Eqs.~\eqref{eq:blendsigma}
and \eqref{eq:blendd} carry an axial insertion depth into the two-dimensional
radial model as an axially averaged absorber. The same insertion axis serves a
second purpose: for every non-control material it is evaluated at the
core-average insertion fraction, because the rod-out/rod-in branch pair also
measures the spectrum shift that inserted absorbers impose on the remainder of
the core.

Whether intermediate depths are \emph{sampled} is then a physical question with
different answers in the two cores, and the blend was measured on the assembled
cores before deciding. In the fast reactor the absorber is not black: the
measured worth curve is close to linear in depth, with $f=0.01$ carrying 1.3\% of
full element worth, so the fast-reactor dataset draws continuous depths and
densely samples the reactivity axis, as real cores hold control elements at
intermediate heights to trim $\keff$. In the pebble bed the same blend is
pathological. B$_4$C in a thermal spectrum is optically black at a few percent
absorber fraction, so a volume-weighted blend makes worth a step function of
depth: $f=0.01$ already carries 53--64\% of full element worth for both the
reflector control rods and the bed shutdown blades. This is the classical
rod-cusping error of homogenized partially inserted absorbers, and it makes an
insertion coordinate that does not mean depth; a first pebble-bed dataset built
on continuous depths was subcritical at every state and was discarded.

Pebble-bed insertion is therefore \emph{binary per element}: each of the ten
reflector control rods and three bed shutdown blades is fully inserted or fully
withdrawn, and the reactivity lever is how many are in and which ones, with the
pattern drawn uniformly so that asymmetric and stuck-element states arise
naturally. The endpoints $f\in\{0,1\}$ are exactly the two transport-measured rod
branches, so the blend of Eq.~\eqref{eq:blendsigma} is never evaluated at an
interior point and the cusping error cannot enter. The cost is that the pebble-bed
reactivity axis is sampled at discrete inserted-element counts rather than
continuously; the benefit is that every rodded state in that database is a state
transport actually measured.

Per-node states are drawn per sample, with the ranges listed in
Table~\ref{tab:sampling}. In the fast reactor a burnup and a temperature are
drawn per assembly, and each control assembly draws its own insertion depth
from the range of its split. In the pebble bed a burnup and a temperature are
drawn per pebble, which is physical rather than cosmetic variability:
multi-pass recirculation means that pebbles of every burnup coexist at every
radius. The number of inserted control rods and of inserted shutdown blades is
drawn from the range of the split and the inserted subset is then chosen
uniformly at random, so asymmetric tilt and stuck-element off-normal states
appear without being imposed. An optional radial burnup
ramp may be blended into the independent draw, with default weight zero, to
represent slower near-wall pebble flow; because its magnitude is proprietary it
is opt-in and the weight used is recorded in each sample. Non-fuel materials
are evaluated at zero burnup. Homogenized materials use the volume- and
flux-weighted transport constants of Section~\ref{sec:groupconstants} -- the
fast-reactor duct blends HT9 and sodium, each pebble node blends its resolved
TRISO and graphite constituents with the surrounding FLiBe, and the barrel,
downcomer, and vessel collapse into a single ring -- so that no homogenization
is performed with representative constants at generation time.

\begin{table}[htbp]
\centering
\caption{Per-node state sampling ranges.}
\label{tab:sampling}
\begin{tabular}{@{}lll@{}}
\toprule
Quantity & Sodium fast reactor & Pebble-bed reactor \\
\midrule
Burnup ($\MWdkg$)        & 2--60, per assembly   & 2--190, per pebble \\
Temperature (K)          & 630--1000, per assembly & 823.15--1100, per pebble \\
Insertion                & Continuous $f\in[0,1]$ & Binary per element \\
Insertion lever          & Depth of each of 13 assemblies & Count and pattern (10 rods, 3 blades) \\
Depletion branch grid ($\MWdkg$) & 0, 2, 20, 40, 60 & 0, 2, 20, 50, 90, 130, 160, 190 \\
Specific power (W/gHM)   & 45                    & 30 \\
Depletion chain          & ENDF/B-VIII.0 fast    & ENDF/B-VIII.0 thermal \\
Group boundary ($G{=}2$) & 0.1 MeV               & 0.625 eV \\
Group boundaries ($G{=}4$) & 9.1188 keV, 0.1 MeV, 1.353 MeV & -- \\
\bottomrule
\end{tabular}
\end{table}

\subsubsection{Finite element solver}
\label{sec:solver}

Following geometry and material assignment, the multigroup neutron diffusion
equations are discretized with the P1 finite element method on the triangular
mesh \citep{duderstadt1976,stacey2007}. Sparse loss and production matrices are assembled from P1 element
stiffness matrices, with the per-group element diffusion coefficient taken as
the arithmetic mean of the three vertex values of the triangle, lumped mass
matrices on the nodal volumes, and Marshak partial-current Robin boundary terms
\citep{marshak1947} added to the stiffness diagonal on each boundary edge, with a
contribution $\alpha L/2$ per endpoint of an edge of length $L$. The Robin
coefficient is written for a general albedo $\beta$, the fraction of the outgoing
partial current returned to the core, $J^- = \beta J^+$:
\begin{equation}
\alpha(\beta) \;=\; \frac{1-\beta}{2\,(1+\beta)} .
\label{eq:albedo}
\end{equation}
Vacuum is thus the $\beta=0$ member of one family, $\alpha = 0.5$, rather than a
separate code path, and $\beta\rightarrow1$ gives $\alpha\rightarrow0$, a
reflective boundary with no leakage. The main datasets are generated at
$\beta=0$; $\beta$ is varied only in the boundary-condition transfer family of
Section~\ref{sec:datafamilies}, where the per-edge coefficient is carried on the
geometry and the nodal value is exposed to the model as an additional node
feature. This yields the
generalized eigenvalue problem of Eq.~\eqref{eq:eigen} in a group-major block
layout of dimension $GN \times GN$. Diagonal blocks carry stiffness, removal, and
the Robin term; off-diagonal blocks carry the in-scatter sources; production
blocks are
\begin{equation}
F_{gg'} \;=\; M\!\left(\chi_g\, \nu\Sigma_{\mathrm{f},g'}\right),
\label{eq:prodblock}
\end{equation}
with $\chi$ taken from the tallied core-average branch spectrum at the state of
that core, falling back to the analytic spectrum of the reactor -- $(0.99,\,0.01)$
for the fast cut and $(1.0,\,0.0)$ for the thermal cut -- if the table carries
none.

Two physical terms enter at this stage and nowhere else. First, axial leakage:
a two-dimensional radial model would otherwise neglect it, so the removal cross
section of each group is augmented with a transverse-buckling term
\citep{stacey2007} as
\begin{equation}
\Sig{r,\mathit{g}} \;\leftarrow\; \Sig{r,\mathit{g}} + D_g \Bz,
\label{eq:buckling}
\end{equation}
with the geometric buckling of the reactor given in Table~\ref{tab:geometry};
the transport models are axially reflective precisely so that this term is not
double counted. Each buckling is $\Bz = (\pi/H_{\mathrm{extrap}})^2$ and is
derived from the axial configuration rather than fitted to a target $\keff$. For
the fast reactor $H_{\mathrm{extrap}} = 135$~cm, an approximately 100~cm active
height plus approximately 15~cm of reflector savings per side. The savings term
is not a refinement: a sodium fast reactor carries an axial reflector and a
sodium plenum above the fuel, and savings are large in a fast core precisely
because $D$ is large, so the flux extrapolates well past the fuel edge. The
corresponding bare-slab value is worth approximately 23\% in $\keff$ and would
place every state of the dataset below critical.

Second, thermal up-scatter for the pebble bed: the tallied
$\Sigma_{\mathrm{s},2\rightarrow1}$ at the temperature of each node adds an
in-scatter source to the fast group at the operator level, keeping the per-node
cross-section row down-scatter-only and the stored schema unchanged. Only the
source term is added there, because the removal side is already present:
$\Sig{r}$ of Eq.~\eqref{eq:removal} is defined as absorption plus out-scatter in
\emph{both} directions, so the thermal-group removal already contains
$\Sigma_{\mathrm{s},2\rightarrow1}$. Adding it a second time at the operator
level is not a small error in this reactor -- for the moderator pebbles the
up-scatter term is roughly 93\% of the thermal removal cross section, and
double counting it drove every pebble-bed state approximately 21{,}000~pcm
subcritical. Up-scatter is an $S(\alpha,\beta)$ thermal-motion effect and
therefore could not defensibly be treated as a temperature-independent constant.

A validated sparse eigensolver -- power iteration on $A^{-1}F$ using a single
sparse LU factorization -- computes the reference flux distribution and
effective multiplication factor for each configuration, converged to
$10^{-9}$ relative in $\keff$ and $10^{-8}$ in the flux $L_2$ norm, with a
sample accepted only if its relative PDE residual falls below $10^{-6}$. The
nodal power density is then computed from the converged flux and the
corresponding fission production cross sections,
\begin{equation}
p_i \;=\; E_{\mathrm{f}} \sum_{g} \nu\Sigma_{\mathrm{f},g}(\mathbf{x}_i)\,\phi_g(\mathbf{x}_i),
\label{eq:power}
\end{equation}
with the energy-per-fission constant folded in. The power head of the model
reuses exactly this relation, so that labels and predictions remain consistent.

\subsubsection{Graph construction}
\label{sec:graph}

After the finite element solution is obtained, each mesh is converted into a
graph representation for graph neural operator training, keeping two graphs
strictly separate. The \emph{physics graph} is the P1 triangulation itself,
stored as elements; it alone defines the assembly of $A$ and $F$ and the
evaluation of the PDE residual. The \emph{message graph} is a distinct
connectivity used only for neural message passing, in which every mesh node
becomes a graph node carrying its spatial coordinates, one-hot material
encoding, boundary indicator and local cross sections.

The message graph is a $k$-nearest-neighbour graph on the mesh nodes, with
$k=10$ by default, built with a KD-tree \todocite{k-d tree, Bentley 1975} and
closed symmetrically, so that it is
stored as a directed \texttt{edge\_index} containing both directions of every
edge and no self-loops. The fixed degree yields clean batching and consistent
GPU utilization on variable-size meshes. Each edge stores physically meaningful interface information:
relative distance and displacement, a material-interface flag, the
harmonic-mean interface diffusion coefficient, and the local differences of the
diffusion coefficient, removal cross section and down-scattering cross section,
with the constants of the fastest group used as the representative interface
scalars. As in graph kernel network formulations of the neural operator
\citep{li2020}, the message graph deliberately admits richer connectivity than
the P1 stencil, allowing information to propagate beyond immediate element neighbors
while preserving local geometric relationships. It is emphasized that $A$ and
$F$ are assembled exclusively from the P1 triangulation; the message graph
exists solely to support neural message passing.

\subsubsection{Database construction}
\label{sec:database}

Each reactor configuration is stored as an independent database sample
containing everything required for supervised and physics-informed training:
node coordinates, triangular element connectivity
$\texttt{elements}[T,3]$, boundary edges $\texttt{boundary\_edges}[B,2]$,
per-node integer material labels, multigroup cross sections, lumped nodal
volumes $\texttt{nodal\_volume}[N]$, the sparse loss and production matrices
$A$ and $F$ of dimension $GN \times GN$, the message-graph connectivity
$\texttt{edge\_index}[2,E]$ with eight-dimensional edge features (distance,
$\Delta x$, $\Delta y$, material-interface flag, harmonic diffusion
coefficient, and the diffusion, removal and down-scattering differences), a
boundary mask, the group fluxes $\texttt{flux}[N,G]$, the effective
multiplication factor $\keff$ and the nodal power density.

Node feature vectors are schema-driven rather than fixed-width. Each vector is
composed of $[x,y]$, a one-hot material block, an ordered cross-section block of
width $D(G) + \Sig{r}(G) + G(G-1)/2$ down-scatter entries $+\ \nuSigf(G)$, and a
boundary flag, with the exact ordering recorded in the sample metadata. Material
is one-hot encoded rather than represented by an ordinal identifier, to avoid
imposing a spurious ordering between materials. At the default $G=2$ this
yields an 18-dimensional feature vector for the eight-material fast reactor and
a 17-dimensional vector for the seven-material pebble bed; the widths follow
directly from the group count and material set and are read from metadata,
never hardcoded. A dataset may append trailing scalar columns after the boundary
flag, which is how the nodal Marshak coefficient of Eq.~\eqref{eq:albedo} is
carried on the boundary-condition transfer family without changing the layout of
any other dataset; the columns present are named in the recorded feature
ordering, and models size their lifting network from it. The resulting layout is
summarized in Table~\ref{tab:features}.

\begin{table}[htbp]
\centering
\caption{Node and edge feature layout at the default group count $G=2$.}
\label{tab:features}
\begin{tabular}{@{}llcc@{}}
\toprule
& Block & Width (general $G$) & Width ($G=2$) \\
\midrule
\multirow{8}{*}{Node}
 & Coordinates $[x,y]$          & 2                & 2 \\
 & One-hot material             & $M$              & 8 (SFR) / 7 (FHR) \\
 & Diffusion coefficient $D_g$  & $G$              & 2 \\
 & Removal $\Sig{r,\mathit{g}}$ & $G$              & 2 \\
 & Down-scatter                 & $G(G-1)/2$       & 1 \\
 & Fission production $\nuSigf$ & $G$              & 2 \\
 & Boundary flag                & 1                & 1 \\
 & Nodal Marshak $\alpha$ (BC family only) & 0 or 1 & 0 or 1 \\
\cmidrule(l){2-4}
 & \textbf{Total}               & $3+M+3G+G(G-1)/2$ & 18 (SFR) / 17 (FHR) \\
\midrule
\multirow{2}{*}{Edge}
 & Distance, $\Delta x$, $\Delta y$, interface flag & 4 & 4 \\
 & Harmonic $D$, $\Delta D$, $\Delta\Sig{r}$, $\Delta\Sigma_{\mathrm{s}}$ & 4 & 4 \\
\cmidrule(l){2-4}
 & \textbf{Total}               & 8                & 8 \\
\bottomrule
\end{tabular}
\end{table}

Geometry metadata describing the reactor type, group structure and fission
spectrum, layout, assembly or pebble-zone positions, per-element control- and
shutdown-insertion states, material counts, and node and edge feature ordering
is retained to enable controlled dataset variability and future scalability
studies. Anything that a sample cannot be reconstructed from otherwise is
recorded there rather than inferred: the mesh refinement level, together with the
distinct levels actually used and whether the mesh is uniform, since the node
count also moves with the reflector ring count and therefore does not identify
the level; the boundary albedo and the resulting per-edge Marshak coefficient;
and the names of any trailing node-feature columns. Every sample additionally
carries an \texttt{xs\_provenance} block
recording the transport code and version, the evaluated data library, the
weighting spectrum, the homogenization and diffusion-coefficient definitions,
the depletion chain and per-zone enrichments, the branch grid, the per-branch
$\keff$, the worst Monte Carlo relative standard deviation in the table, the
interpolation rules and a converged flag. This block is never absent, since the
material modules raise an error at import if their table is missing, so a
sample cannot exist without transport provenance. Pebble-bed samples further
carry an explicit \texttt{modeling\_assumptions} block recording where the
model knowingly departs from the licensed methodology -- an independent rather
than flow-path-correlated burnup field, a slice of the cylindrical section only,
a solid rather than channeled graphite reflector, binary rather than partial
element insertion, and multigroup diffusion in place of full-core
three-dimensional Monte Carlo with burnup -- so that those departures are
auditable from the dataset alone.

\subsubsection{Verification and sample validation}
\label{sec:verification}

Generating group constants from transport makes the labels traceable, but it
does not by itself demonstrate that multigroup diffusion reproduces transport,
which is the claim on which the dataset rests. A dedicated verification harness
closes that gap by solving identical core states in both ways: with the P1 FEM
diffusion model used to label the dataset, and with a full-core
continuous-energy OpenMC eigenvalue calculation. It reports, per state, the
difference $\Delta k = k_{\mathrm{diff}} - k_{\mathrm{CE}}$, the reactivity bias
in pcm, the control element worth from both codes, and the root-mean-square and
maximum relative deviation between the diffusion power density and the OpenMC
fission-rate distribution binned in common radial rings. Because the transport
model is axially reflective, each state is solved twice with the diffusion
model -- once at $\Bz=0$ to match the transport model, which is the pair from
which the bias is measured, and once with the buckling of the dataset -- and the
reactivity difference between the two is reported as an explicit axial debit, so
that the axial treatment remains visible rather than hidden inside the bias.
States are run at uniform burnup, temperature, and insertion, since the per-node
scatter of the dataset has no single continuous-energy counterpart.

The same harness carries the group-resolution sensitivity study, by re-solving
the fast-reactor states on the four-group table of
Section~\ref{sec:groupconstants} while reusing the continuous-energy statepoints
unchanged, the reference being independent of the group structure. This
separates the two candidate explanations for a reactivity bias in the labels:
too coarse a group collapse, or the diffusion approximation itself. The result
(Section~\ref{sec:res-verification}) attributes the residual to the latter, which
is why the datasets are generated at $G=2$ and the four-group structure is kept
as a sensitivity study rather than a second dataset.

Every generated sample is additionally checked before it reaches training.
The checks comprise schema-derived tensor shapes and feature widths, one-hot
validity, finiteness, sparse operator dimensions $[GN,GN]$, a non-empty boundary
mask, valid triangles and boundary edges with in-range indices, strictly
positive lumped nodal volumes summing to the meshed core area (partition of
unity), presence of the defining structural and fuel materials in the mesh,
absence of self-loops in the message graph, a relative PDE residual below
$10^{-6}$ for the reference solution, and a physically plausible $\keff$.

Training, validation and test splits are kept disjoint on control insertion,
which is the axis that most strongly moves the solution. In the fast reactor the
insertion depth ranges are 0--0.5, 0.5--0.75 and 0.75--1.0; in the pebble bed,
where insertion is binary, the splits are disjoint on the \emph{number} of
inserted elements, at zero to five control rods and zero or one shutdown blade
for training, six or seven and two for validation, and eight to ten and three
for test. The remaining axes are variability rather than separation and are
allowed to overlap: the fast-reactor enrichment-boundary ring and reflector-ring
counts, and the pebble packing seed, moderator-pebble fraction and per-node
burnup and temperature, all vary within each split.

This is a deliberately hard split, and its consequence must be stated plainly:
$\keff$ tracks insertion, so the validation and test sets occupy $\keff$ ranges
the model never sees in training and the $\keff$ head is required to
\emph{extrapolate}, not to interpolate. In the pebble-bed database the realized
ranges are 0.986--1.076 for training, 0.929--0.990 for validation and
0.833--0.905 for test. Reported $\keff$ errors on validation and test should be
read in that light; a split that shuffled configurations at random would report
far smaller errors while measuring far less.

\subsubsection{Dataset families}
\label{sec:datafamilies}

Two main databases are generated, one per benchmark, each of 7000 samples split
5000/1000/1000 by the plans above, at fast-reactor mesh level $s=0$ and vacuum
boundaries. Every sample converged, with a worst relative PDE residual of
approximately $1.2\times10^{-8}$ across both databases. Typical sizes are 3421
nodes for the fast reactor at 217--271 assemblies and 7098 nodes for the pebble
bed at approximately 1730 pebbles.

Two further families vary a single generation axis in order to pose
generalization questions that the main databases cannot. Both are fast-reactor
only.

\paragraph{Mesh resolution} The refinement level $s$ of
Section~\ref{sec:geometry} is swept over levels 1--3, giving approximately 4700,
13\,800 and 28\,200 nodes per core, together with a non-uniform level in which
the inner and outer assemblies carry different levels. The family is
\emph{paired}: every level solves the same physical cores -- the same insertion
depths, the same per-assembly burnup and temperature, the same ring counts --
because the number of random draws depends only on the assembly count, which
refinement does not change, so re-seeding replays an identical draw sequence at
each level. The pairing is asserted at generation time rather than assumed.
Without it, a difference across levels would confound discretization with
configuration, since insertion alone moves $\keff$ by far more than refinement
does. The non-uniform level is the sharper probe of the two aggregation rules
compared in Section~\ref{sec:kernel}: uniform refinement changes node density
everywhere at once, which a message-passing model can partly absorb by
rescaling, whereas a mixed mesh makes density vary \emph{within} a single graph,
so the core no longer has one local node count for an unweighted sum to track.

\paragraph{Boundary condition} The albedo $\beta$ of Eq.~\eqref{eq:albedo} is
held disjoint across splits at $[0,0.40)$, $[0.40,0.55)$ and $[0.55,0.80]$, the
test range extrapolating, while insertion, ring counts, burnup and temperature
are drawn identically in all three, so that a failure to generalize is
attributable to the boundary condition and not to a core state the model never
saw. This family drops the outer shield ring, and that departure from the
production layout is not optional. The shield is optically thick, so on the
production core the boundary is very nearly invisible: sweeping $\beta$ from 0 to
0.8 moves $\keff$ by 0.5~pcm and the normalized flux by 0.065\%. Removing the
shield raises those to 1445~pcm and 11.08\%, and removing the radial reflector as
well to 5371~pcm and 29.15\%. The reflector-only core is used, keeping the
radial reflector a fast core physically requires while making the outer boundary
a live part of the problem; on the production geometry every model would appear
to generalize perfectly by ignoring the boundary condition entirely, and the
result would be an artifact. The choice is recorded in every sample of the
family. The pebble bed admits no equivalent: at the published vessel stack the
same sweep moves $\keff$ by 1.4~pcm, and thinning that stack would depart from
the dimensional specification the benchmark is pinned to, so boundary-condition
transfer is claimed for the fast reactor only.

%--------------------------------------------------------------
\subsection{Physics-informed graph neural operator framework}
\label{sec:pigno}

\subsubsection{Feature lifting}
\label{sec:lifting}

Prior to message passing, node and edge features are independently projected
into higher-dimensional latent spaces using multilayer perceptrons with SiLU
activation functions \todocite{SiLU activation},
\begin{equation}
\mathbf{h}_i^{(0)} = \mathrm{MLP}_{\mathrm{node}}(\mathbf{x}_i),
\qquad
\mathbf{e}_{ij} = \mathrm{MLP}_{\mathrm{edge}}(\mathbf{a}_{ij}),
\label{eq:lift}
\end{equation}
where $\mathbf{x}_i$ is the node feature vector of Table~\ref{tab:features} and
$\mathbf{a}_{ij}$ the corresponding edge feature vector. These lifting networks
transform the raw physical reactor quantities into latent representations
better suited to graph operator learning while preserving the underlying
reactor physics.

\subsubsection{Kernel integration layers}
\label{sec:kernel}

The lifted node representations are processed through multiple graph neural
operator kernel integration layers \citep{li2020,kovachki2023}. Within each layer, information from
neighboring nodes is aggregated by residual message passing,
\begin{equation}
\mathbf{h}_i^{(\ell+1)} = \mathbf{h}_i^{(\ell)}
+ \mathrm{LN}\!\left(
\mathcal{A}_{j\in\mathcal{N}(i)}\!\left[\,
m_{ij}^{(\ell)} \,\right]
\right),
\qquad
m_{ij}^{(\ell)} =
\kappa_\theta\!\left(\mathbf{h}_i^{(\ell)},\,\mathbf{h}_j^{(\ell)},\,\mathbf{e}_{ij}\right),
\label{eq:kernel}
\end{equation}
where each message depends on the latent representations of the neighboring
nodes together with their associated edge features describing geometric and
material interactions, $\kappa_\theta$ is the learned kernel and $\mathrm{LN}$
denotes layer normalization \todocite{layer normalization, Ba et al. 2016}.
Layer normalization is applied after each
message-passing block to improve optimization stability and to facilitate
training of deeper graph operator architectures.

The aggregation $\mathcal{A}$ is where a graph neural network and a graph
neural operator part company, and it is left selectable so that the difference
can be measured rather than asserted. Three rules are implemented,
\begin{equation}
\mathcal{A}^{\mathrm{sum}} = \sum_{j\in\mathcal{N}(i)} m_{ij},
\qquad
\mathcal{A}^{\mathrm{vol}} =
\frac{\sum_{j\in\mathcal{N}(i)} V_j\, m_{ij}}{\sum_{j\in\mathcal{N}(i)} V_j},
\qquad
\mathcal{A}^{\mathrm{raw}} = \sum_{j\in\mathcal{N}(i)} V_j\, m_{ij},
\label{eq:aggregation}
\end{equation}
with $V_j$ the lumped nodal volume of Section~\ref{sec:solver}. The unweighted
scatter-add $\mathcal{A}^{\mathrm{sum}}$ is a graph neural network aggregation;
it carries no quadrature weight and its value therefore depends on how many
neighbours the mesh happens to supply. The volume-weighted average
$\mathcal{A}^{\mathrm{vol}}$ is the Nystr\"om approximation of the kernel
integral
\begin{equation}
(\mathcal{K}\phi)(x_i) = \int_\Omega
\kappa(x_i, y)\,\phi(y)\,\mathrm{d}y
\;\approx\; \sum_{j\in\mathcal{N}(i)} V_j\,\kappa(x_i, y_j)\,\phi_j ,
\label{eq:nystrom}
\end{equation}
in which the weight belongs to the \emph{source} node $j$ because the integral
is taken over $y$; this is the default and is the form the operator claim
rests on. The unnormalized $\mathcal{A}^{\mathrm{raw}}$ is the strict Nystr\"om
estimate and is retained to show why the normalization is needed: with a
fixed-degree $k$-nearest-neighbour graph the neighbourhood shrinks as the mesh
refines, so the raw sum loses magnitude with resolution while the normalized
form does not. Section~\ref{sec:res-discretization} reports all three, which is
what attributes any operator-versus-network difference to the volume quadrature
rather than to the name of the model.

The limitation of that construction is stated here rather than left for a
reader to find. Discretization invariance in the graph-neural-operator sense
requires both a fixed physical integration domain and quadrature weights that
converge to the underlying measure. Volume weighting supplies the second;
it does not supply the first, since a fixed-degree neighbourhood contracts as
$N$ grows. What is claimed and measured in this work is therefore not the
limiting theorem but the part of it that acts over the two- to fourfold
refinement range the study spans.

\subsubsection{Prediction heads}
\label{sec:heads}

Following message passing, separate prediction heads generate the physical
quantities of interest. Node-level prediction layers estimate the neutron flux
in each energy group, while graph-level pooling followed by a multilayer
perceptron predicts the multiplication factor,
\begin{equation}
\hat{\phi}_{i,g} = \mathrm{MLP}_{\phi}\!\left(\mathbf{h}_i^{(L)}\right),
\qquad
\hat{k}_{\mathrm{eff}} = \mathrm{MLP}_{k}\!\left(
\frac{1}{N}\sum_{i=1}^{N}\mathbf{h}_i^{(L)}\right).
\label{eq:heads}
\end{equation}
Rather than introducing an additional neural prediction head, the nodal power
density is computed analytically from the predicted flux using
Eq.~\eqref{eq:power}, thereby improving physical consistency and reducing the
number of learned quantities.

\subsubsection{Physics-informed objective}
\label{sec:loss}

Physics-informed constraints are incorporated directly into the optimisation
objective, in the manner of physics-informed learning
\citep{raissi2019,sarkar2025}, by minimising the residual of the governing
discrete diffusion eigenvalue problem,
\begin{equation}
\mathcal{R} \;=\; A\,\hat{\boldsymbol{\phi}}
- \frac{1}{\hat{k}_{\mathrm{eff}}}\,F\,\hat{\boldsymbol{\phi}},
\label{eq:residual}
\end{equation}
which encourages the predicted solutions to satisfy the underlying reactor
physics. The complete loss function comprises four components,
\begin{equation}
\mathcal{L} \;=\;
\lambda_{\phi}\,\mathcal{L}_{\phi}
+ \lambda_{k}\,\mathcal{L}_{k}
+ \lambda_{\mathrm{PDE}}\,\mathcal{L}_{\mathrm{PDE}}
+ \lambda_{\mathrm{BC}}\,\mathcal{L}_{\mathrm{BC}},
\label{eq:loss}
\end{equation}
namely the supervised neutron flux loss, the multiplication factor loss, the
PDE residual loss and the boundary condition loss, with $\lambda$ the
corresponding weights. The flux and $\keff$ losses are evaluated in normalized
units for scale stability, whereas the residual of Eq.~\eqref{eq:residual} is
evaluated in physical units, since $A$ and $F$ are physical operators.

The boundary term requires a caveat, because the boundary condition is enforced
twice over and only one of the two routes is active by default. The Marshak
condition of Section~\ref{sec:solver} is a Robin condition assembled into $A$
itself, so it is already carried by $\mathcal{L}_{\mathrm{PDE}}$, and the
reference boundary flux it produces is deliberately \emph{nonzero}. A
zero-flux penalty $\mathcal{L}_{\mathrm{BC}} = \langle
\hat{\phi}^{2}\rangle_{\partial\Omega}$ therefore does not express the boundary
condition of this discretization; it contradicts the reference solution, and
driving it to zero would drive the prediction away from the labels. The term is
consequently weighted at $\lambda_{\mathrm{BC}} = 0$ for every result reported
here, and is retained in Eq.~\eqref{eq:loss} for two purposes: as the reported
boundary residual diagnostic of Section~\ref{sec:res-physics}, and as an active
penalty for a Dirichlet-zero discretization, where it would be the correct
statement of the condition. The physics constraint in this work is thus the
operator residual, with the boundary condition entering through the assembled
operator rather than as an independent penalty.

%--------------------------------------------------------------
\subsection{Benchmark models}
\label{sec:baselines}

The proposed PI-GNO framework is evaluated against three baselines, each
representing a distinct and widely used answer to the same problem: a
grid-based neural operator, a non-graph neural operator, and a mesh-based graph
neural network. They are listed in Table~\ref{tab:baselines}.

The Fourier neural operator and MeshGraphNet baselines are taken from NVIDIA
PhysicsNeMo 2.1.1 \todocite{NVIDIA PhysicsNeMo 2.1.1} rather than reimplemented,
so that the comparison against the proposed model is not self-refereed. The
deep operator network \todocite{DeepONet, Lu et al. 2021} is implemented in this
work, following the original branch--trunk--inner-product formulation, because
it is not part of the PhysicsNeMo core library; DeepONet has previously been
applied to two-dimensional neutron transport \citep{kobayashi2024}.

\begin{table}[htbp]
\centering
\caption{Benchmarked architectures and the capability each is used to probe.
``Irregular geometry'' means the model consumes the unstructured mesh without
resampling; ``discretization transfer'' means the same trained weights may be
evaluated at a different mesh resolution; ``new boundary condition'' means the
model has an explicit channel through which an unseen boundary condition can
enter at test time.}
\label{tab:baselines}
\begin{tabular}{@{}llccc@{}}
\toprule
Model & Source & Irregular & Discretization & New \\
      &        & geometry  & transfer       & BC \\
\midrule
FNO          & PhysicsNeMo 2.1.1 & --      & \checkmark & -- \\
DeepONet     & this work         & partial & partial    & -- \\
MeshGraphNet & PhysicsNeMo 2.1.1 & \checkmark & --      & -- \\
PI-GNO       & this work         & \checkmark & \checkmark & \checkmark \\
\bottomrule
\end{tabular}
\end{table}

One frequent statement about the Fourier neural operator is not made here,
because it is false: FNO learns its kernel in Fourier space and is genuinely
resolution-flexible, which is its central claim. Its limitation in this setting
is narrower and specific. FNO requires a \emph{uniform Cartesian grid}, and
neither a hexagonal duct lattice nor a randomly packed pebble bed is one, so the
input field must be rasterized onto a grid and the prediction interpolated back
to the mesh nodes before it can be scored. That round trip loses information
regardless of how well the network fits, and the loss is a property of the
geometry rather than of the architecture. It is measured directly, without any
trained model, in Section~\ref{sec:res-floors}, and reported beside every FNO
result, so that ``FNO performed worse on irregular geometry'' decomposes into
the part attributable to the grid and the part attributable to the model.
Similarly, DeepONet is marked partial rather than absent on both axes: its trunk
evaluates at arbitrary coordinates and therefore at the mesh nodes directly,
incurring no output interpolation error, but its branch reads a fixed sensor
lattice on a fixed domain, so its input representation is as grid-bound as
FNO's.

A second qualification matters as much. PI-GNO with the unweighted aggregation
$\mathcal{A}^{\mathrm{sum}}$ of Eq.~\eqref{eq:aggregation} \emph{is} a graph
neural network in exactly the respect under test; the operator behaviour comes
from the volume quadrature, not from the model's name. The aggregation
ablation of Section~\ref{sec:res-discretization} is therefore presented with
$\mathcal{A}^{\mathrm{sum}}$ as the network-side control rather than as a third
variant of the proposed model.

\subsubsection{Architectural fairness}
\label{sec:fairness-arch}

Four choices remove the most obvious confounders from the comparison.

\begin{enumerate}
\item MeshGraphNet receives the \emph{same} $k$-nearest-neighbour message graph
and the same eight-dimensional edge features as PI-GNO, not the finite element
triangulation. Had the two models seen different graphs, any accuracy
difference would have been attributable to the graph rather than to the
processor.
\item MeshGraphNet's decoder emits a latent vector to which the \emph{same}
flux and $\keff$ heads of Eq.~\eqref{eq:heads} are attached, so the comparison
is processor against processor with the heads held fixed.
\item Every model, including the baselines, computes nodal power from its
predicted flux through the same parameter-free relation of
Eq.~\eqref{eq:power}. Power error therefore reflects flux error alone, and no
model can win on power by fitting the power labels with a free head.
\item FNO predicts on its grid and is interpolated back to the mesh nodes
\emph{before} being scored. Scoring it on its own grid would conceal precisely
the error the baseline exists to expose. Rasterization onto the grid is
nearest-neighbour, since cross sections are piecewise constant per material and
blending them across an interface would invent materials that do not exist;
interpolation from the grid back to the mesh is bilinear, the flux being smooth
and the linear choice being the more generous one for FNO.
\end{enumerate}

%--------------------------------------------------------------
\subsection{Training and evaluation methodology}
\label{sec:training}

All models are trained on identical training, validation and test datasets,
with the split construction of Section~\ref{sec:datafamilies}. Validation data
are used for model selection and for the one hyperparameter that is tuned per
model; the test set contains reactor configurations outside the training range
of the separating axis, so every reported error is an extrapolation error.

\subsubsection{Controls on the comparison}
\label{sec:fairness-train}

Five controls are applied, each answering an objection that would otherwise
render the comparison unfalsifiable.

\paragraph{A common parameter budget} Every model is sized to
$4.0\times10^{5}$ trainable parameters by bisecting a single width knob, so that
``the proposed model was simply larger'' has a recorded answer. Width rather
than depth is the knob, because depth changes the receptive field, which is
itself part of what is under study for the graph models. The achieved counts are
given in Table~\ref{tab:budget}. Three of the four models land within
0.6\% of the budget. FNO lands 8.0\% \emph{above} it: its spectral convolution
parameters scale with the square of the latent channel count, so the budget is
not exactly reachable, and the nearest attainable configuration was taken on
the side that favours the baseline rather than starves it. Actual counts, not
the nominal budget, are reported.

\begin{table}[htbp]
\centering
\caption{Parameter-budget matching. One width knob per model is bisected to a
common target of $4.0\times10^{5}$ trainable parameters.}
\label{tab:budget}
\begin{tabular}{@{}llrrr@{}}
\toprule
Model & Width knob & Value & Parameters & Deviation \\
\midrule
PI-GNO       & latent dimension     & 25 & 398\,605 & $-0.3\%$ \\
MeshGraphNet & hidden dimension     & 75 & 399\,765 & $-0.1\%$ \\
DeepONet     & branch hidden width  & 35 & 402\,384 & $+0.6\%$ \\
FNO          & latent channels      & 10 & 432\,199 & $+8.0\%$ \\
\bottomrule
\end{tabular}
\end{table}

\paragraph{A common learning-rate grid} Each model selects its learning rate
from the same grid $\{3\times10^{-4},\,10^{-3},\,3\times10^{-3}\}$ on the
validation split, using the same shortened budget. This is the only per-model
tuning performed, so ``the baselines were undertuned'' is answered by equal
treatment rather than by unequal effort.

\paragraph{One training loop} All four models pass through a single training
and evaluation implementation, with one normalization procedure and one metric
function. Optimization is AdamW \todocite{AdamW, Loshchilov and Hutter 2019}
with weight decay $10^{-5}$, gradient-norm clipping at 1.0, a five-epoch linear
warm-up followed by cosine annealing over a 200-epoch budget, batch size four,
and early stopping after 30 epochs without validation improvement. Evaluation
runs at batch size one so that per-core and per-material metrics remain
per core. Normalization statistics are fitted on the training split only, are
persisted alongside each checkpoint, and are applied unchanged to validation,
test and transfer data, so that a checkpoint is always evaluated with the
transform it was trained under.

\paragraph{A data-only accuracy comparison} The architecture comparison is run
with $\lambda_{\mathrm{PDE}}=0$ for \emph{every} model, PI-GNO included. This is
the single most important control in the study: were PI-GNO to retain its
physics loss while the baselines had none, any advantage would be attributable
to the objective rather than to the architecture. The physics term is
reintroduced separately, and credited separately, in the ablation of
Section~\ref{sec:res-physics}.

\paragraph{Repetition} Every cell is run with three random seeds, and results
are reported as mean $\pm$ standard deviation across seeds.

\subsubsection{Experiment matrix}
\label{sec:matrix}

The five experiments of Table~\ref{tab:experiments} expand into 35 cells, each
a model paired with one sweep value, and 105 training runs at three seeds per
cell. E5 is a reporting pass over timings the training runs already record and
adds no runs of its own.

\begin{table}[htbp]
\centering
\caption{Experiment matrix. ``Cells'' counts model $\times$ sweep-value
combinations; each cell is run at three seeds.}
\label{tab:experiments}
\begin{tabular}{@{}llcl@{}}
\toprule
Experiment & Database & Cells & Question \\
\midrule
E1 (fast reactor) & main SFR      & 4 & architecture accuracy, data-only \\
E1 (pebble bed)   & main FHR      & 4 & the same, thermal spectrum, random packing \\
E2 (fast reactor) & main SFR      & 6 & physics-loss weight, graph models \\
E2 (pebble bed)   & main FHR      & 6 & the same \\
E3 transfer       & mesh family   & 4 & discretization transfer, all models \\
E3 aggregation    & mesh family   & 3 & operator-versus-network mechanism \\
E4 transfer       & albedo family & 4 & boundary-condition transfer \\
E4 with physics   & albedo family & 4 & does the physics term aid BC transfer \\
\bottomrule
\end{tabular}
\end{table}

E2 and the aggregation ablation are restricted to the graph models, which are
the only ones that consume the assembled operators $A$ and $F$ directly. E3 and
E4 are run on the fast reactor only, for reasons that are physical rather than
practical and are stated rather than omitted. The pebble-bed mesh carries one
node per pebble, so its node count \emph{is} the reactor and cannot be refined
without changing the problem; and its boundary condition is unmeasurable at the
published vessel stack, as quantified in Section~\ref{sec:datafamilies}.

E3 trains at the middle uniform refinement level and evaluates the same
checkpoints at every other level, so that both a coarsening and a refinement are
tested; training at either end of the family would test only one direction.
Transfer evaluations reuse the trained checkpoint rather than retraining per
target, retraining being a different question.

\subsubsection{Metrics}
\label{sec:metrics}

The same quantities are recorded per core for every experiment and every model.
Numerical accuracy is quantified by the relative $L_2$ error,
\begin{equation}
\varepsilon_{L_2} =
\frac{\lVert \hat{\boldsymbol{\phi}} - \boldsymbol{\phi} \rVert_2}
     {\lVert \boldsymbol{\phi} \rVert_2},
\label{eq:l2}
\end{equation}
evaluated per energy group, together with the mean absolute error, the
root-mean-square error, the coefficient of determination $R^2$, and the error
in the multiplication factor, reported in pcm. The same relative $L_2$ measure
is applied to the nodal power density and, restricted to the nodes of each
material, to the per-material breakdown that locates where the error lives.
Physical consistency is assessed by the root-mean-square residual of
Eq.~\eqref{eq:residual} and by the boundary flux diagnostic of
Section~\ref{sec:loss}. Computational cost is recorded as inference wall-clock
time per core alongside the node count, for comparison against the reference
solver timings of Section~\ref{sec:res-cost}.

\subsubsection{Measuring discretization transfer}
\label{sec:res-transfer-metric}

Discretization transfer requires more care than the other axes, because the
obvious metric measures the wrong thing. Scoring a model against \emph{each
mesh's own} finite element reference does not measure discretization
invariance, since those references are not the same field: they differ from one
another by the finite element convergence error, quantified in
Section~\ref{sec:res-floors}. Scored that way, a model is charged for the
discretization error of its label, and a model that fit every mesh perfectly
would appear maximally mesh-\emph{dependent}. Two metrics are therefore
reported together.

\paragraph{Invariance, with no reference involved} The same trained model is
evaluated on two meshes of the same physical core, one prediction is
interpolated onto the other's nodes, and the disagreement between them is
measured. A discretization-invariant operator scores zero on this metric
irrespective of whether it is accurate.

\paragraph{Accuracy against a fixed fine mesh} All predictions are additionally
scored against one finely resolved reference solution, with each mesh's own
distance from that reference --- its discretization floor --- reported beside
it, since no model can reasonably be faulted for error below the floor of the
data it was given.

Both are necessary, because invariance alone is gameable: a constant output is
perfectly invariant, and in a preliminary run an undertrained model scored
better invariance than the finite element references themselves while being
uniformly wrong at every resolution.

A scale convention is required for either metric to mean anything. The
eigensolver of Section~\ref{sec:solver} returns an eigenvector normalized to
unit $L_2$ norm over its $GN$ entries, so the magnitude of a reference flux
vector is an artifact of the node count, falling as $N^{-1/2}$: the maximum
fast-group value across the mesh family of Section~\ref{sec:datafamilies} is
0.0375, 0.0212 and 0.0147 at successive refinement levels. All cross-mesh
comparisons are therefore renormalized by $\int \sum_g \phi_g \,\mathrm{d}V$,
one scalar applied to all groups, so that the physical spectral ratio between
groups survives the renormalization.

%=========================================================================
\section{Results and discussion}
\label{sec:results}

Results are presented in two parts. Sections~\ref{sec:res-verification} and
\ref{sec:res-floors} report properties of the databases and of the reference
solver that hold independently of any trained model, and that bound what any
model in the study can achieve. Sections~\ref{sec:res-accuracy} to
\ref{sec:res-cost} report the trained models against those bounds.

Three qualifications apply to every number that follows and are stated once
here.

\begin{enumerate}
\item \emph{All reported errors are extrapolation errors.} The splits of
Section~\ref{sec:datafamilies} are disjoint on control insertion, so the
validation and test sets occupy $\keff$ ranges and rod configurations the models
never see in training. These are not independent and identically distributed
holdouts, and the absolute errors are correspondingly larger than a shuffled
split would report. The comparison \emph{between} models, at a common split, is
the part of the result that carries meaning.
\item \emph{Cross-section provenance.} Every sample carries its transport
provenance: the Monte Carlo code and version, the evaluated data library, the
branch grid, the per-branch $\keff$ and its Monte Carlo uncertainty. No
hand-tuned or representative cross-section library exists anywhere in the
pipeline.
\item \emph{The unweighted aggregation is a graph neural network.} In the
aggregation ablation of Section~\ref{sec:res-discretization},
$\mathcal{A}^{\mathrm{sum}}$ is presented as the network-side control, not as a
third flavour of the proposed operator.
\end{enumerate}

\subsection{Dataset and reference solution verification}
\label{sec:res-verification}
% TODO: table of diffusion vs. continuous-energy transport per verification
% state (dk, bias in pcm, control worth in both codes, axial debit, radial
% power RMS/max), plus the four-group sensitivity study of Sec. 2.4.

\subsection{Error floors independent of any model}
\label{sec:res-floors}

\subsubsection{Finite element discretization floor}
\label{sec:res-femfloor}

The reference labels inherit the convergence error of the P1 discretization
that produced them. Refining the fast-reactor mesh over levels 3 to 6 gives
successive $\keff$ gaps of 541, 329 and 221~pcm, a ratio of 1.49 against the
1.50 that first-order convergence predicts, so the observed order is
$p \approx 1.0$. Richardson extrapolation on that sequence places even the
finest mesh used here, at $N \approx 1.1\times10^{5}$ nodes, approximately
1100~pcm from the continuum limit, and the production mesh of the main
fast-reactor database approximately 4700~pcm from level 3.

This bounds what the labels are, and it is why they are described as
self-consistent rather than as converged. Within one refinement level the
labels are internally consistent to the eigensolver tolerance, with a worst
relative residual of order $10^{-8}$, which is what a machine learning
benchmark requires: every model is fit to, and scored against, the same
discrete operator. It is not what a comparison against continuous-energy
transport requires, and any reactivity bias reported in
Section~\ref{sec:res-verification} must be read with this floor beside it.

\subsubsection{Rasterization floor for the grid baseline}
\label{sec:res-rasterfloor}

The mesh-to-grid-to-mesh round trip that the Fourier neural operator requires
loses information before any network is trained. Pushing the \emph{reference}
flux through that round trip therefore bounds the FNO error from below.
Table~\ref{tab:rasterfloor} reports the resulting floor for both benchmarks at
three grid resolutions.

\begin{table}[htbp]
\centering
\caption{Rasterization floor: relative $L_2$ error incurred by resampling the
reference flux onto a uniform grid and interpolating it back to the mesh nodes.
No trained model is involved. Every FNO result is bounded below by these
values.}
\label{tab:rasterfloor}
\begin{tabular}{@{}lcccc@{}}
\toprule
& \multicolumn{2}{c}{Fast reactor} & \multicolumn{2}{c}{Pebble bed} \\
\cmidrule(lr){2-3}\cmidrule(lr){4-5}
Grid & group 1 & group 2 & group 1 & group 2 \\
\midrule
$64^2$  & 2.7\% & 2.4\% & 3.7\% & 9.1\% \\
$128^2$ & 1.2\% & 1.1\% & 1.8\% & 4.7\% \\
$192^2$ & 0.7\% & 0.6\% & 1.1\% & 3.1\% \\
\bottomrule
\end{tabular}
\end{table}

The thermal group of the pebble bed is the sharp case. A $128^2$ grid spanning
the 420~cm core diameter resolves 3.3~cm per cell against a 4~cm pebble, so the
grid cannot represent the pebble-scale thermal flux structure at all, and no
FNO trained on that grid can fall below 4.7\% thermal flux error regardless of
its fit. FNO is accordingly given a grid sweep over $\{64^2, 96^2, 128^2,
192^2\}$ and reported at its best configuration, with the corresponding floor
alongside.

\subsubsection{Discretization floor of the mesh family}
\label{sec:res-meshfloor}

Table~\ref{tab:probefloor} reports the two reference quantities defined in
Section~\ref{sec:res-transfer-metric}, measured over 30 paired configurations
of the mesh family --- the same physical cores meshed six ways. The invariance
column is the finite element solution's \emph{own} mesh dependence, and is the
number a model must beat in order to be more discretization-invariant than the
discretization it was trained on. The floor column is each mesh's distance from
the finest reference.

\begin{table}[htbp]
\centering
\caption{Reference behaviour of the paired mesh family, 30 configurations.
Invariance is measured against level 3 and the floor against the level-6 truth
mesh; both are relative $L_2$ errors after the common renormalization of
Section~\ref{sec:res-transfer-metric}. These are properties of the finite
element solutions alone.}
\label{tab:probefloor}
\begin{tabular}{@{}llrcccc@{}}
\toprule
& & & \multicolumn{2}{c}{Invariance vs.\ level 3}
    & \multicolumn{2}{c}{Floor vs.\ level 6} \\
\cmidrule(lr){4-5}\cmidrule(lr){6-7}
Mesh & Character & $N$ & group 1 & group 2 & group 1 & group 2 \\
\midrule
Level 0 & different meshing rule & 3\,821  & 0.0963 & 0.0584 & 0.1131 & 0.0742 \\
Level 1 & coarser                & 5\,275  & 0.0655 & 0.0667 & 0.0815 & 0.0827 \\
Mixed   & density varies within one graph & 9\,271 & 0.0863 & 0.0570 & 0.0977 & 0.0716 \\
Level 2 & training discretization & 15\,447 & 0.0158 & 0.0172 & 0.0315 & 0.0342 \\
Level 3 & finer                  & 31\,432 & ---    & ---    & 0.0155 & 0.0172 \\
Level 6 & truth mesh             & 114\,265 & ---   & ---    & ---    & ---    \\
\bottomrule
\end{tabular}
\end{table}

Two features of Table~\ref{tab:probefloor} are worth stating explicitly.

First, the references disagree with one another substantially. The fast-group
flux differs by 6.5\% between levels 1 and 3 after renormalization, and $\keff$
moves 3059~pcm between levels 1 and 2 on the \emph{same physical core}. This is
the quantitative reason a per-mesh reference cannot serve as the transfer
metric.

Second, the mixed mesh has approximately 1.8 times the node count of level 1
yet a worse floor, and the cause is measured rather than assumed. It is not
element quality: the triangle aspect-ratio distributions of the two meshes are
indistinguishable, with a 99th percentile of 13.47 and a maximum of 30.16. It
is domain area. The Delaunay gap-filling filter of Section~\ref{sec:geometry}
rejects more triangles where a finely meshed assembly abuts a coarsely meshed
one, so the mixed mesh represents 0.94\% less core area than the truth mesh
against 0.76\% for level 1, and the floor tracks the area deficit across the
whole family. The correct reading is that a mesh with more nodes covered less
of the core, not that more nodes gave worse physics.

\subsection{Flux and multiplication factor accuracy}
\label{sec:res-accuracy}
% TODO (E1): table of relative L2 per group, MAE, RMSE, R^2, k_eff error in
% pcm and power relative L2, for all four models on both benchmarks, at
% lambda_pde = 0, mean +- std over three seeds. The FNO row carries its
% rasterization floor from Table~\ref{tab:rasterfloor}.
% TODO: per-material breakdown table.
% TODO: figure, predicted vs. reference flux maps and pointwise error.

\subsection{Physical consistency and the physics-informed objective}
\label{sec:res-physics}
% TODO (E2): PDE residual and boundary diagnostic vs. lambda_pde in
% {0, 0.1, 1.0} for PI-GNO and MeshGraphNet, both benchmarks. This is where the
% physics-informed claim of the title is earned or is not.

\subsection{Discretization transfer}
\label{sec:res-discretization}
% TODO (E3): model columns for Table~\ref{tab:probefloor} -- invariance and
% accuracy vs. the level-6 truth mesh, for all four models trained at level 2.
% TODO (E3 aggregation): PI-GNO under A_sum, A_vol and A_raw of
% Eq.~\eqref{eq:aggregation}, presented with A_sum as the GNN-side control, on
% the uniform levels and on the mixed mesh.

\subsection{Boundary-condition transfer}
\label{sec:res-bc}

The albedo family of Section~\ref{sec:datafamilies} was regenerated without the
outer shield ring because the boundary condition is otherwise invisible. On the
generated reflector-only database the albedo axis moves $\keff$ by
approximately 1410~pcm end to end, against 0.5~pcm on the production layout, so
the experiment measures a boundary condition that is actually present in the
solution.

One caveat belongs with the result rather than after it. Sweeping the albedo is
a \emph{synthetic} axis: exactly one value of $\beta$ reproduces the solution
of the full radial geometry, and the remainder are boundary conditions that no
physical reactor presents. The experiment is a legitimate test of
generalization across a family of boundary conditions, and it is not a physical
parameter sweep.

% TODO (E4): error vs. held-out beta for all four models; and (E4 with physics)
% the same for the graph models at lambda_pde in {0, 0.1}, which asks whether
% the operator residual --- known exactly at a new boundary condition, without
% labels --- aids transfer.

\subsection{Computational cost}
\label{sec:res-cost}

Table~\ref{tab:solvercost} gives the reference solver cost that any inference
time must be compared against. Assembly of $A$ and $F$ accounts for 1.5--2\% of
the total; the sparse LU factorization dominates. Wall-clock time scales as
$N^{1.27}$ to $N^{1.31}$ across the family, consistent with nested-dissection
sparse LU in two dimensions. All timings are single-threaded CPU.

\begin{table}[htbp]
\centering
\caption{Reference solver cost per core, by database, on CPU. $N$ and the times
are means over the timed samples of each split. Assembly is the construction of
$A$ and $F$; the solve is the sparse LU factorization with power iteration to a
relative residual below $10^{-8}$.}
\label{tab:solvercost}
\begin{tabular}{@{}lrrrr@{}}
\toprule
Database & $N$ & Assemble (s) & Solve (s) & Total (s) \\
\midrule
Fast reactor, main   & 4\,279   & 0.005 & 0.318  & 0.323  \\
Pebble bed, main     & 7\,093   & 0.012 & 0.748  & 0.760  \\
Mesh family, level 2 & 15\,562  & 0.019 & 1.699  & 1.718  \\
Mesh family, level 3 & 31\,666  & 0.048 & 4.301  & 4.348  \\
Mesh family, level 6 & 115\,114 & 0.138 & 22.142 & 22.280 \\
\bottomrule
\end{tabular}
\end{table}

These numbers constrain what may honestly be claimed. A fast-reactor core
solves in 0.32~s, so generating the entire 7000-sample database cost
approximately 38 minutes of CPU time, and ``faster than the reference solver''
is not by itself a compelling argument for a single core. The defensible
arguments are the many-query setting --- design sweeps, uncertainty
quantification and optimization, where $10^4$ to $10^6$ solves are required ---
and the large-mesh setting, where the reference takes 22~s per core and scales
superlinearly. A ratio between a CPU solver and a GPU model is also reported
only with that label attached, since such a ratio measures hardware as much as
method.

% TODO (E5): inference time per core and parameter count per model, and the
% accuracy-versus-inference-time Pareto front against the solver times above.

\subsection{Generalisation beyond the training distribution}
\label{sec:res-generalisation}
% TODO: figure, error vs. insertion depth across the disjoint split ranges;
% error vs. reflector-ring count and enrichment-boundary ring (fast reactor);
% error vs. moderator pebble fraction (pebble bed).
% OPTIONAL, cheap once E4 exists: test-time physics refinement. At an unseen
% boundary condition A is known exactly and requires no labels, so PI-GNO can
% take a few gradient steps on the PDE residual of the test sample itself. No
% purely data-driven baseline has that channel available.

%=========================================================================
\section{Conclusions}
\label{sec:conclusions}

% TODO: write after results. Cover, and only if the data support it:
%   - PI-GNO accuracy vs. the FEM reference in flux, k_eff and power (E1)
%   - contribution of the PDE-constrained objective (E2, lambda_pde ablation)
%   - discretization transfer, and whether the volume quadrature rather than
%     the graph is what produced it (E3 + aggregation ablation)
%   - boundary-condition transfer, with the synthetic-axis caveat (E4)
%   - inference cost in the many-query setting, not as a per-core speedup (E5)
%   - future work (3D cores, thermal-hydraulic feedback, transients)
% Each claim must be stated against the floor it is measured over: the FNO
% comparison against Table~\ref{tab:rasterfloor}, the transfer results against
% Table~\ref{tab:probefloor}. Do not claim FNO cannot handle varying
% resolution -- the claim is that it requires a uniform Cartesian grid.

%=========================================================================
\section*{CRediT authorship contribution statement}
% TODO: agree the split with the co-author before submission. CRediT taxonomy:
% Conceptualization, Methodology, Software, Validation, Formal analysis,
% Investigation, Resources, Data curation, Writing -- original draft,
% Writing -- review & editing, Visualization, Supervision, Project
% administration, Funding acquisition.
%
% \textbf{Author Name:} TODO.

\section*{Declaration of competing interest}
The authors declare that they have no known competing financial interests or
personal relationships that could have appeared to influence the work reported
in this paper.

\section*{Declaration of generative AI and AI-assisted technologies in the writing process}
% Required by Elsevier when applicable; delete if not applicable.

\section*{Data availability}
% State one of: data will be made available on request; a repository link; or
% that the code and generated database are openly available.

\section*{Acknowledgments}
% Funding sources and computing allocations.

%=========================================================================
\bibliographystyle{elsarticle-harv}   % author-year, as used by Ann. Nucl. Energy
\bibliography{refs}

\end{document}
